\section{Auditoría de fuentes: el scaling no es una curva, sino todo un Operating System}

Esta clase la imparte Ben Mann, cofundador de Anthropic, en febrero de 2025. La sesión va de GPT-2/GPT-3 y las scaling laws a las distributed failure del frontier training, los loss spike, RLHF/Constitutional AI, la evaluación de seguridad, la Responsible Scaling Policy (RSP), la interpretabilidad mecanicista y las distintas disciplinas de lanzamiento de los productos de chat y de la API. El repositorio conserva la portada oficial y 1,025 subtítulos con marca de tiempo; esta nota verifica los mecanismos con artículos y materiales oficiales de Anthropic, y no proyecta hacia la clase el estado posterior de los modelos, los ingresos, la RSP ni los productos.

\begin{warningbox}{Lo que dijo la clase y los límites de versión}
Cifras como «crecimiento de diez veces en el último año» o «el coding segment creció diez veces en tres meses» se toman solo como la descripción del workload que dio el ponente en ese momento, no como conclusiones financieras auditadas. La clase usa el relato de los AI Safety Level (ASL) de la era 2025; al 11 de agosto de 2026, Anthropic ya había publicado el 24 de febrero de 2026 la RSP v3.0, que adopta un marco actualizado con capability thresholds, required safeguards, entre otros elementos. Esta nota primero explica fielmente la clase histórica y después señala por separado la evolución posterior.
\end{warningbox}

\begin{importantbox}{Los cuatro ciclos cerrados de un frontier model program}
\begin{enumerate}
  \item \textbf{Scaling loop}: hypothesis $\rightarrow$ pilot runs $\rightarrow$ fit/forecast $\rightarrow$ compute allocation;
  \item \textbf{Training loop}: telemetry $\rightarrow$ anomaly $\rightarrow$ checkpoint/replay $\rightarrow$ fix/resume;
  \item \textbf{Safety loop}: preentrenamiento $\rightarrow$ post-training $\rightarrow$ evaluation $\rightarrow$ safeguards/deployment;
  \item \textbf{Product loop}: chat experiment $\rightarrow$ behavior evidence $\rightarrow$ API contract $\rightarrow$ migration/deprecation.
\end{enumerate}
\end{importantbox}

\begin{knowledgebox}{Términos clave: hay que distinguir cuatro tipos de «escala»}
\textbf{Training scale} se refiere a los parámetros, los token y los FLOPs de entrenamiento; \textbf{serving scale}, al volumen de solicitudes, el context, la latency y la accelerator fleet; \textbf{capability scale}, al éxito en las tareas, la generalización y el tool use; \textbf{risk scale}, al impacto potencial que surge de la combinación de la capacidad del modelo, la forma de acceso y las safeguard failure. El crecimiento de los ingresos o del tráfico no demuestra por sí mismo la capability, y el avance en los benchmark tampoco demuestra automáticamente la safety.
\end{knowledgebox}

\subsection{Por qué el crecimiento cambia los supuestos de ingeniería}

Un crecimiento rápido de la demanda no significa solo «comprar más GPU». Training e inference compiten por la capacity, los usuarios de la API exigen un contract estable, los usuarios del chat esperan iteraciones más rápidas, el equipo de seguridad tiene que atender más adversarial traffic y el on-call enfrenta un costo de falla más alto. El program necesita colocar capability, compute, reliability y safeguards en un mismo modelo de recursos, en lugar de optimizarlos por separado.

\lecturefigure{01-scaling-program.png}{El frontier scaling es un program en el que están acoplados demand, capability, compute y reliability.}{Subtítulos locales 00:00--01:00; diagrama conceptual redibujado.}

\begin{knowledgebox}{Lectura de la figura: el crecimiento genera retroalimentación positiva, pero también riesgos que se acumulan}
La demand aporta ingresos y tareas reales; la capability eleva la adoption; el compute sostiene el entrenamiento y el serving; la reliability mantiene la confianza de los usuarios. Cualquier eslabón desequilibrado amplifica el efecto en sentido inverso: un capability release sin serving capacity provoca colas, el traffic sin abuse control amplía el riesgo, el entrenamiento que desplaza a la inference perjudica a los clientes actuales y una API regression convierte la actualización del modelo en un outage del negocio.
\end{knowledgebox}

\subsection{Resumen de la sección}

El objeto de un scaling program es el ciclo de vida completo. Las secciones siguientes no tomarán un «modelo más grande» como respuesta única, sino que examinarán cómo la predicción, el entrenamiento, la recuperación, la alineación, la evaluación, el lanzamiento y la compatibilidad se sostienen en conjunto.
